% Pursuing Stacks (cotes 134-2 to 134-8): archive cote 134-4, site page 44 (PDF page 45), typescript page 132.
% typescript. Draft compared: ps3.tex:2223-2283.
% AI-transcribed by Claude Opus 5.5 (claude-opus-5-5, adaptive thinking, effort high), Message Batches API, batch pilot (msgbatch_0189Nhf1xB615YWnztubtLfc), 2026-11-12; not yet checked by eye.
\scanpage{134-4/44}
\mspage{132}

The first statement is just \xout{partxof} the "homotopy lemma" part b)
(page 99), the second follows by the argument sketched above.

We're about back now \xout{\unsure{xtxthexpoint}} to the context we started with
three days ago (par. 48 , page 98 and following), where we started with
a W (viewed as a notion of ``weak equivalence''), and in terms of W con-%
structed various homotopy notions - namely those, we would now say,
corresponding to the homotopy structure defined by the set of all intervals \margin{\unsure{$\mathbb{I}$} (with a brace along these lines)}
in M which are W-aspheric \xout{xxxx} (i.e. I \inserted{is} W-aspheric over e\hand{*}). As a matter of \margin{* 11.4. by which we mean that $I\to e$ is ``universally in $W$''. The terminology used here for ``$W$-aspheric'' is highly ambiguous, cf. \deleted{\unsure{dist}} discussion p. 181 and following}
fact, we were a little stricter still, by restricting to intervals
which are \inserted{moreover} ``disjoint'' (and which we called``homotopy intervals''\xout{\unsure{for}} relative
to W ), but this restriction now \inserted{definitely} appears as awkward and artificial. I
will henceforth call \uline{homotopy interval\xout{\unsure{s}}} (with respect to W), any interval
(not necessarily a separated one) which is W-aspheric. Let $h_W$ be the
corresponding homotopy structure on M , which is a weak homotopy interval
structure admitting the set of all W-homotopy intervals as a generating
set of weak homotopy intervals. \hand{(}\unsure{Clearly}, there will be \inserted{many} weak homotopy
intervals for this structure, which are far from being W-aspheric i.e.
far from being homotopy intervals.\hand{)} Of course, \xout{the} as stated in the pre-%
ceding proposition, $h_W$ and W are \inserted{strictly} compatible i.e.
\[ W_{(h_W)} \hand{$\subset$} W \]
i.e. any $h_W$-homotopism is in W (\inserted{i.e. is} a``weak equivalence''). As a matter of fact,
the definition of homotopy notions in terms of W we gave in loc. cit.
were just the widest one we could think of by that time, which would
ensure the ``compatibility'' of these notions with W , in a sense which
was'nt technically clear (not even definable at that point) as it is now,
but however reasonably clear in\xout{x} terms of \inserted{mathematical} ``bon sens''. At present though
the question arises rather naturally whether the homotopy structure $h_W$
we selected ``au flair'' by that time is indeed the best one, namely the
widest one, we could get. More explicitly, this means whether the
homotopy structure $h_W$ is the \uline{widest} (in terms of \deleted{the} \inserted{natural} order relation
considered in the previous paragraph) among all those which are\xout{x} compatib-%
le with W in the strict sense $W_{\hand{h}} \hand{$\subset$} W$ . Now this is certainly not so , if
we are not a little more specific about \xout{which} restricting to homotopy
structures definable in terms of a weak homotopy interval structure.
For instance, if we \inserted[typed]{take for W a \xout{\unsure{weak}} homotopism structure on M ,} \xout{\unsure{takexforxWxxxx(M,W)xxthexmodelizerxofxallxKanxcom-}}
\xout{\unsure{plexesxandxweakxequivalencesxcorrespondingxto}} compatible with products,
corresponding to a homotopy structure h , to say $h_W$ is the ``best'' would
\xout{xxxxxx} imply that W \inserted{itself} can be described in terms of weak homotopy intervals

\margin{44}
